

\section{Le choc de la démocratisation (1975) : Le grec devient une option-refuge}

L'histoire du système éducatif français de la seconde moitié du XXe siècle est dominée par une
tension dialectique fondamentale entre une volonté politique d'unification structurelle — incarnée
par la figure du « collège unique » — et la persistance, voire la recomposition, de logiques de
différenciation sociale et scolaire. Le point I.3.\cite{I-3-1-ref1} de notre plan d'étude, intitulé
« Le choc de la démocratisation (Réforme Haby 1975) : le 'Collège unique' et le grec comme option-
refuge pour éviter la carte scolaire », nous invite à explorer l'une des manifestations les plus
subtiles et les plus symboliques de cette tension : l'instrumentalisation des enseignements de
Langues et Cultures de l'Antiquité (LCA), et singulièrement du grec ancien, dans les stratégies de
reproduction sociale et d'évitement de la mixité scolaire.

La loi du 11 juillet 1975, portée par le ministre René Haby sous la présidence de Valéry Giscard
d'Estaing, constitue l'acte de naissance législatif du collège unique. En mettant fin à la
ségrégation institutionnelle des filières qui prévalait jusqu'alors au sein du premier cycle du
second degré, cette réforme ambitionnait de réaliser la promesse républicaine de l'égalité des
chances en offrant à toute une classe d'âge un savoir commun, retardant ainsi les processus de
sélection et d'orientation.\cite{I-3-1-ref1} Pourtant, l'analyse sociologique et historique démontre
que l'unification administrative n'a pas aboli la hiérarchisation sociale des parcours ; elle l'a
déplacée. Face à l'hétérogénéité nouvelle des publics scolaires, perçue comme un risque de
déclassement ou de nivellement par le bas par les catégories sociales favorisées, l'option grec
ancien a progressivement acquis une fonction latente : celle de servir d'outil de distinction et de
contournement de la carte scolaire.

Ce rapport de recherche se propose d'analyser de manière exhaustive les mécanismes par lesquels le
grec ancien, discipline académique par excellence, est devenu une « option-refuge » stratégique.
Pour ce faire, nous mobiliserons un cadre théorique croisant la sociologie de l'éducation (notamment
les travaux de Pierre Bourdieu, Marie Duru-Bellat et Pierre Merle), l'histoire des politiques
éducatives et l'analyse des flux scolaires basée sur les données de la DEPP. Nous examinerons
successivement la genèse du collège unique comme rupture historique, le fonctionnement de la carte
scolaire et ses dispositifs dérogatoires, la sociologie spécifique de l'option grec par rapport au
latin, et enfin les conséquences de ces stratégies sur la ségrégation scolaire interne et externe.



\subsection{Genèse et tensions du collège unique : du « choc » haby à la démocratisation ségrégative}

Pour saisir la portée stratégique du grec ancien après 1975, il est impératif de replacer la réforme
Haby dans la longue durée de l'histoire scolaire française. Le « choc » ressenti par les acteurs du
système éducatif ne provient pas uniquement de la réforme elle-même, mais de la rupture brutale
qu'elle opère avec un modèle séculaire de ségrégation des publics.



\subsection{L'héritage de la segmentation : l'école d'avant 1975}

Avant la mise en œuvre de la réforme Haby, le système éducatif français était structuré par une
logique d'ordres distincts, héritée du XIXe siècle et de la distinction entre l'école du peuple (le
primaire et le primaire supérieur) et l'école des élites (le secondaire, les lycées et collèges).
Comme le rappellent les historiens de l'éducation, cette organisation résultait d'une conception où
deux réseaux de scolarisation coexistaient sans se mélanger, conformes aux visions de Guizot ou de
Destutt de Tracy, plutôt qu'à l'idéal unificateur de Condorcet.\cite{I-3-1-ref2}

Jusqu'au début des années 1960, et malgré la réforme Berthoin de 1959 qui prolongeait la scolarité
obligatoire jusqu'à 16 ans, le premier cycle du second degré restait profondément cloisonné. La
création des Collèges d'Enseignement Secondaire (CES) en 1963 par la réforme Fouchet-Capelle avait
certes initié un regroupement physique des élèves (« l'école sous un même toit »), mais elle
maintenait une ségrégation pédagogique stricte à travers trois filières étanches :

1. \textbf{La filière I (Type Lycée)} : Elle dispensait un enseignement classique ou moderne long,
assuré par des professeurs certifiés ou agrégés (professeurs de lycée). C'était la voie royale
menant au baccalauréat, fréquentée majoritairement par les enfants de la bourgeoisie et des classes
moyennes supérieures. Les langues anciennes y occupaient une place centrale.

2. \textbf{La filière II (Type CEG)} : Héritière des Cours Complémentaires, elle proposait un
enseignement moderne court, assuré par des instituteurs spécialisés ou des PEGC (Professeurs
d'Enseignement Général de Collège). Elle accueillait une population issue des classes moyennes et
populaires en ascension sociale.

3. \textbf{La filière III (Type Transition)} : Elle était destinée aux élèves en difficulté
scolaire, issus quasi-exclusivement des classes populaires, et débouchait rapidement sur la vie
active ou l'apprentissage.\cite{I-3-1-ref3}

Dans cette architecture, le grec ancien n'avait pas besoin de jouer un rôle d'option-refuge : la
sélection s'opérait en amont, par l'affectation dans la filière I. L'étude du grec était alors une
composante \og naturelle\fg{} du curriculum des élites, un marqueur d'appartenance à la section
classique, protégé par les murs invisibles mais infranchissables des filières.



\subsection{La rupture haby : l'idéal de l'indifférenciation et ses conséquences}

La loi du 11 juillet 1975 constitue un tournant majeur en supprimant ces filières. René Haby,
recteur et ministre, porte le projet d'unifier le corps social par l'école. Le « collège unique »
impose un cursus commun à tous les élèves de la 6ème à la 3ème. L'objectif est double : élever le
niveau général de formation pour répondre aux besoins d'une économie moderne (la « massification »)
et assurer la justice sociale en retardant le moment de l'orientation.\cite{I-3-1-ref1}

Cependant, cette unification structurelle a engendré un « choc » sociologique et pédagogique. En
mélangeant dans les mêmes classes des élèves aux capitaux culturels et aux niveaux scolaires très
disparates, la réforme a créé une situation d'hétérogénéité inédite. Pour les familles favorisées,
habituées à l'entre-soi des filières nobles (Lycées et CES filière I), le collège unique a été perçu
comme une menace de déclassement et de nivellement par le bas.\cite{I-3-1-ref3} L'historien Claude
Lelièvre et le sociologue Christian Nique soulignent bien que cette réforme, bien qu'initiée par un
gouvernement de droite libérale (Giscard d'Estaing), a été vécue comme une rupture radicale avec la
tradition de l'élitisme républicain.\cite{I-3-1-ref6}

C'est dans ce contexte de désarroi des élites et des classes moyennes intellectuelles que les
options facultatives, et particulièrement le grec ancien, vont changer de fonction. De disciplines
de culture, elles vont devenir des instruments de défense. Comme l'institution ne garantit plus la
séparation des publics par les filières, les usagers vont chercher à rétablir cette séparation par
le jeu des options. C'est le passage d'une ségrégation institutionnelle subie à une ségrégation
stratégique choisie.



\subsection{Le concept de « démocratisation ségrégative »}

Pour théoriser ce phénomène, il convient de mobiliser le concept de « démocratisation ségrégative »
développé par le sociologue Pierre Merle. Ce concept décrit une situation paradoxale où l'ouverture
de l'accès à l'école (démocratisation quantitative ou massification) s'accompagne d'un maintien,
voire d'un renforcement, des écarts sociaux dans l'accès aux filières les plus
valorisées.\cite{I-3-1-ref8}

Pierre Merle distingue trois types de démocratisation :

\begin{itemize}\item \textbf{La démocratisation égalisatrice} : Réduction des écarts sociaux de réussite et d'orientation.\item \textbf{La démocratisation uniforme} : Décalage général des scolarités vers le haut sans modification des écarts relatifs entre groupes sociaux.\item \textbf{La démocratisation ségrégative} : Accès généralisé à un niveau d'enseignement (ici, le collège et la classe de 3ème) qui s'accompagne d'une différenciation qualitative accrue des parcours. Les enfants des classes populaires accèdent au collège unique \og standard\fg{}, tandis que les enfants des classes favorisées construisent des parcours protégés via les options distinctives (Grec, Latin, Allemand LV1).\cite{I-3-1-ref8}\end{itemize}Dans le cadre du collège post-1975, le grec ancien devient l'un des vecteurs principaux de cette
ségrégation. Il permet de recréer, au sein d'un établissement formellement unifié, des filières
invisibles qui reproduisent la stratification sociale antérieure.



\subsection{L'ingénierie de l'évitement : la carte scolaire et le « parcours scolaire particulier »}

L'utilisation du grec ancien comme option-refuge ne relève pas seulement d'une préférence culturelle
; elle s'inscrit dans une mécanique administrative précise liée à la sectorisation, communément
appelée « carte scolaire ».



\subsection{La carte scolaire : un carcan administratif devenu enjeu social}

Mise en place progressivement à partir de 1963 pour gérer les flux démographiques, la carte scolaire
assigne chaque élève à un établissement public en fonction de son lieu de résidence. Avec
l'avènement du collège unique, elle est devenue l'instrument principal de la mixité sociale, mais
aussi la cible de toutes les stratégies d'évitement.\cite{I-3-1-ref5}

La rigidité théorique de la carte scolaire a toujours été tempérée par un système de dérogations.
Ces dérogations, gérées par les Directions des Services Départementaux de l'Éducation Nationale
(DSDEN), sont accordées selon des motifs hiérarchisés. Si les motifs médicaux ou sociaux (bourses)
sont prioritaires, un motif pédagogique s'est rapidement imposé comme la clé de voûte des stratégies
d'évitement des classes moyennes et supérieures : le motif de « parcours scolaire particulier
».\cite{I-3-1-ref5}



\subsection{Le grec ancien comme motif légitime de dérogation}

Le « parcours scolaire particulier » permet à une famille de demander l'inscription de son enfant
dans un collège hors secteur si cet établissement propose un enseignement qui n'existe pas dans le
collège de rattachement. Or, contrairement à l'anglais ou à l'espagnol, présents partout, le grec
ancien est une « option rare », dont l'offre est inégalement répartie sur le territoire. Il est
souvent concentré dans les collèges de centre-ville historiquement prestigieux ou dans les
établissements qui cherchent à maintenir une mixité sociale par le haut.\cite{I-3-1-ref16}

Cette rareté fait du grec un levier puissant. Une famille résidant dans le secteur d'un collège jugé
« difficile » ou « moyen » peut légitimement demander une dérogation pour inscrire son enfant dans
un collège plus réputé, au motif que l'enfant souhaite ardemment débuter l'apprentissage du grec
ancien (généralement en 4ème ou 3ème selon les époques). L'administration, ne pouvant juger de la
sincérité de la vocation helléniste, est contrainte d'accepter la demande si des places sont
disponibles.\cite{I-3-1-ref5}

Les textes officiels, notamment les circulaires sur la préparation de la rentrée et les notes de
service sur l'affectation (comme celles liées à \textit{Affelnet} plus récemment), mentionnent
explicitement ce critère. Le décret n° 2005-1309 et les circulaires subséquentes (2007, 2008) ont
institutionalisé cette pratique, même si l'assouplissement de 2007 visait officiellement à donner
plus de poids aux critères sociaux.\cite{I-3-1-ref5} Le grec ancien apparaît ainsi dans les
documents administratifs comme une variable d'ajustement des flux, permettant une mobilité choisie
pour une minorité initiée.



\subsection{L'assouplissement de la carte scolaire : accélérateur de la concurrence}

L'histoire de la carte scolaire est marquée par des vagues d'assouplissement qui ont,
paradoxalement, renforcé le rôle stratégique des options. Dès 1984, sous Alain Savary, puis en 1987
sous René Monory, des expérimentations de « libre choix » ont été menées.\cite{I-3-1-ref19} Ces
politiques ont montré que les familles les plus dotées en capital culturel utilisaient massivement
les marges de manœuvre offertes pour fuir les établissements stigmatisés.

L'assouplissement généralisé de 2007, sous la présidence de Nicolas Sarkozy et le ministère de
Xavier Darcos, a marqué un point de non-retour. En supprimant officiellement la rigidité de la carte
scolaire (tout en maintenant des priorités), l'État a mis les collèges en concurrence directe. Pour
survivre et éviter la fuite vers le privé, les chefs d'établissement des collèges publics ont dû
développer des « produits d'appel ». Maintenir une section de grec ancien, même avec un très faible
effectif, est devenu vital pour conserver les élèves des classes moyennes et supérieures du
secteur.\cite{I-3-1-ref21} Le grec est devenu un outil de marketing scolaire, un signal de qualité
envoyé aux parents consommateurs d'école.



\subsection{Sociologie de la distinction : pourquoi le grec et pas le latin?}

Si les langues anciennes en général servent de refuge, le grec ancien occupe une place spécifique,
plus sélective et plus distinctive que le latin. L'analyse bourdieusienne de la distinction permet
de comprendre cette hiérarchie.



\subsection{Latin vs grec : distinction de masse et distinction d'élite}

Le latin, débutant généralement en 5ème, a connu une certaine forme de démocratisation, ou du moins
de massification relative au sein des classes moyennes. Il est perçu comme une aide au français, une
base culturelle utile. En revanche, le grec ancien, débutant plus tard (4ème ou 3ème) et souvent en
cumul avec le latin, opère une « sur-sélection ».\cite{I-3-1-ref13}

Les statistiques de la DEPP sont éloquentes. Alors que le latin conserve une base de recrutement
sociologique relativement large (bien que biaisée), le grec est devenu l'apanage quasi-exclusif des
enfants de cadres supérieurs et d'enseignants.

\begin{itemize}\item \textbf{Données sociales comparatives} : Les élèves d'origine sociale défavorisée sont quasiment absents des cours de grec ancien. Selon les données de la DEPP (RERS 2024 et 2025), si environ 15\% des élèves défavorisés peuvent étudier le latin (souvent perçu comme utile pour le français), le pourcentage tombe à des niveaux négligeables pour le grec (\<1\%).\cite{I-3-1-ref24}\item \textbf{La stratégie du cumul} : L'élève helléniste est presque toujours un latiniste. Il cumule deux options lourdes, signalant une capacité de travail et une conformité scolaire exceptionnelles. Ce cumul agit comme un double filtre social.\end{itemize}

\subsection{La valeur symbolique de l'inutile : bourdieu et les humanités}

Pierre Bourdieu, dans \textit{La Distinction} et \textit{La Noblesse d'État}, analyse comment la
culture classique sert à légitimer la domination sociale. Le grec ancien, par son caractère «
gratuit », « désintéressé » et apparemment « inutile » sur le marché du travail, est le marqueur
absolu de la distinction bourgeoise.\cite{I-3-1-ref27}

L'étude du grec fonctionne comme une « épreuve » initiatique qui consacre l'héritier. Elle prouve
que l'élève dispose du temps libre et de la distance au monde nécessaire pour se plonger dans des
savoirs antiques, loin des urgences matérielles. Bourdieu parle de l'efficacité « magique » de cet
enseignement : peu importe que les élèves sachent réellement lire couramment le grec à la fin (ce
qui est rarement le cas), l'important est qu'ils l'aient fait. Le passage par la classe de grec
marque l'habitus, inculque une posture, une aisance avec la culture légitime qui sera valorisée dans
les grands oraux des concours (ENA, ENS).\cite{I-3-1-ref29}



\subsection{Stratégies familiales et rôle des mères}

Les travaux de Marie Duru-Bellat et d'Agnès van Zanten sur les stratégies familiales montrent que ce
sont souvent les mères, gestionnaires du parcours scolaire de l'enfant, qui identifient et activent
ces leviers.\cite{I-3-1-ref33} Le choix du grec n'est pas laissé au hasard ou au seul goût de
l'enfant ; il est le fruit d'un calcul rationnel (ou d'un sens pratique ajusté) visant à placer
l'enfant dans la « bonne » classe ou le « bon » établissement.

De plus, une dimension de genre apparaît : les filles sont statistiquement plus nombreuses à choisir
les options littéraires et les langues anciennes que les garçons, anticipant des carrières où les
humanités restent valorisées ou acceptant plus volontiers la docilité scolaire requise par
l'apprentissage \og grammaire-traduction\fg{}.\cite{I-3-1-ref34} Cependant, pour les garçons des
élites, le maintien du grec (souvent couplé aux mathématiques d'excellence) reste une voie royale
vers les classes préparatoires les plus prestigieuses.



\subsection{La fabrication des « classes de niveau » (ségrégation interne)}

L'efficacité de l'option-refuge ne se limite pas au changement d'établissement (évitement externe) ;
elle est redoutable pour contourner la mixité au sein même du collège (évitement interne). C'est ce
que l'on appelle la constitution des « classes de niveau » ou des « classes protégées ».



\subsection{La gestion de l'hétérogénéité par les chefs d'établissement}

Face à l'hétérogénéité du collège unique, les chefs d'établissement sont soumis à une double
contrainte : assurer la réussite de tous (mission de service public) et retenir les élèves
performants des milieux favorisés (impératif de compétitivité). Pour résoudre cette équation
impossible, ils utilisent les options comme des outils de regroupement.\cite{I-3-1-ref4}

Concrètement, l'administration regroupe les élèves hellénistes (souvent peu nombreux) avec d'autres
élèves ayant choisi des options sélectives (Section Européenne, Bilangue allemand, Classes à
Horaires Aménagés Musique \- CHAM). Cela crée mécaniquement une classe d'élite, où se concentrent
les meilleurs élèves et où les problèmes de discipline sont moindres.

\begin{itemize}\item \textbf{L'effet de pair} : Dans ces classes, l'émulation joue à plein. Les enseignants, souvent les plus expérimentés ou les plus exigeants, y dispensent un cours plus dense, plus rapide.\cite{I-3-1-ref38}\item \textbf{L'effet de protection} : Les parents le savent parfaitement. Inscrire son enfant en grec, c'est lui garantir qu'il ne sera pas dans la classe qui concentre les élèves en grande difficulté ou perturbateurs. C'est une assurance contre le risque scolaire.\end{itemize}

\subsection{Le grec ancien : une garantie supérieure au latin}

Si le latin permet déjà ce type de regroupement dès la 5ème, le grec en 3ème (ou 4ème) permet
d'affiner la sélection à l'approche du lycée. Il agit comme un filtre final avant l'orientation en
Seconde. Dans certains collèges, la « classe de grec » est explicitement identifiée par tous
(élèves, profs, parents) comme la « bonne classe », créant un sentiment de relégation chez les
autres élèves exclus de ce cercle vertueux.\cite{I-3-1-ref11}

Cette pratique est dénoncée par les sociologues comme une forme de « ségrégation intra-établissement
» qui vide le collège unique de sa substance, mais elle est tolérée, voire encouragée tacitement par
l'institution, car elle permet de maintenir la paix sociale et de garder les classes moyennes dans
le secteur public.\cite{I-3-1-ref5}



\subsection{Évolution statistique et ruptures historiques (1975-2024)}

L'analyse des données quantitatives sur le long terme permet de valider l'hypothèse d'une discipline
devenue refuge pour une minorité, plutôt qu'une discipline démocratisée.



\subsection{Analyse des effectifs : l'érosion continue}

Les données issues des \textit{Repères et Références Statistiques} (RERS) de la DEPP montrent une
chute drastique des effectifs, qui contraste avec l'explosion démographique du collège.

\begin{table}[h!]\small\centering\begin{tabular}{| p{2cm} | p{2.2cm} | p{2.2cm} | p{2.2cm} |}\hline\textbf{Année Scolaire} & \textbf{Niveau 4ème (\%)} & \textbf{Niveau 3ème (\%)} & \textbf{Contexte Historique et Politique} \\ \hline\textbf{1975-1976} & \~3-4 \% (est.) & \~3-4 \% (est.) & Mise en place de la réforme Haby. \\ \hline\textbf{1984-1985} & 1,8 \% & 2,0 \% & Premiers assouplissements carte scolaire (Savary). \\ \hline\textbf{1995-1996} & 2,0 \% & 2,0 \% & Stabilité relative avant la réforme Bayrou.\cite{I-3-1-ref24} \\ \hline\textbf{1998-1999} & \textit{N/A} & 1,8 \% & Effet de la réforme 1996 (report en 3ème). \\ \hline\textbf{2005-2006} & 0,1 \% & 1,5 \% & Marginalisation confirmée.\cite{I-3-1-ref25} \\ \hline\textbf{2019-2020} & 0,1 \% & 1,8 \% & Réforme du collège 2016 (EPI) et retouches Blanquer.\cite{I-3-1-ref26} \\ \hline\end{tabular}\end{table}Note de lecture : Les chiffres indiquent le pourcentage d'élèves d'une classe d'âge suivant
l'option. Sources : DEPP, RERS.\cite{I-3-1-ref24}

On observe que le grec, qui concernait une frange significative des élèves de la filière I avant
1975, est devenu une discipline confidentielle (autour de 1,5\% à 2\% des élèves de 3ème). Cette
contraction numérique a renforcé son caractère distinctif : ce qui est rare est cher (socialement).



\subsection{La rupture de 1996 : le report en 3ème}

Un moment clé de cette histoire est la réforme menée par François Bayrou en 1996. Devant la baisse
des moyens et la volonté de promouvoir d'autres enseignements, l'initiation au grec ancien a été
repoussée de la 4ème à la 3ème.\cite{I-3-1-ref40}

Cette décision, justifiée par des arguments budgétaires et pédagogiques (concentrer l'effort sur une
année), a eu un effet pervers majeur : la « rupture de série » statistique.

\begin{itemize}\item \textbf{Effondrement de la base} : En supprimant la 4ème, on a coupé le vivier de recrutement.\item \textbf{Sur-sélection} : Ne choisissent plus le grec en 3ème que les élèves ultra-motivés, ayant déjà réussi leur parcours de 6ème-5ème-4ème, et quasi-systématiquement déjà latinistes. Le grec est devenu une option de finition pour l'excellence, déconnectée d'une logique de découverte pour tous.\item \textbf{Conséquence sur la carte scolaire} : Le motif de dérogation \og grec ancien\fg{} est devenu encore plus efficace, car l'offre s'est raréfiée dans les collèges modestes, obligeant les familles à viser les lycées-collèges de centre-ville qui maintenaient l'option.\cite{I-3-1-ref40}\end{itemize}

\subsection{Les réformes récentes : de l'epi au « choc des savoirs »}

Plus récemment, la réforme du collège de 2016 (Najat Vallaud-Belkacem) a tenté de diluer les langues
anciennes dans des Enseignements Pratiques Interdisciplinaires (EPI) « Langues et Cultures de
l'Antiquité », provoquant une levée de boucliers des défenseurs des humanités qui y voyaient la fin
de l'enseignement linguistique rigoureux.\cite{I-3-1-ref45}

Jean-Michel Blanquer est revenu en partie sur ces dispositions en 2018, rétablissant une option
facultative plus traditionnelle, avant que le « Choc des Savoirs » annoncé ne pose de nouvelles
questions sur les groupes de niveaux.\cite{I-3-1-ref1} À chaque étape, la question de l'option grec
cristallise les débats sur l'élitisme : faut-il le sauver pour préserver l'excellence (vision
conservatrice) ou le supprimer pour garantir l'égalité (vision égalitariste)? La réalité est que,
quelle que soit la réforme, les familles favorisées trouvent le moyen d'utiliser ce qui reste de
l'option pour se distinguer.



\subsection{La barrière pédagogique : méthodes et reproduction}

Enfin, il ne faut pas négliger la dimension proprement pédagogique de cette ségrégation. Le grec
n'est pas seulement une case dans l'emploi du temps ; c'est un contenu, une méthode.



\subsection{La méthode « grammaire-traduction » comme filtre}

L'enseignement du grec ancien en France est resté longtemps, et reste en partie, fidèle à la méthode
traditionnelle dite « grammaire-traduction ». Cette approche met l'accent sur la mémorisation
intensive de la morphologie (déclinaisons, verbes irréguliers) et la gymnastique intellectuelle de
la version.\cite{I-3-1-ref50}

Cette méthode favorise intrinsèquement les élèves disposant d'un \textit{habitus} scolaire conforme
aux attentes bourgeoises : docilité, capacité d'abstraction, rapport scriptural au langage. Comme le
souligne Bourdieu, le formalisme de cet enseignement, déconnecté de l'usage vivant de la langue,
fonctionne comme un test d'appartenance sociale.\cite{I-3-1-ref27} Pour les élèves de milieux
populaires, l'aridité de l'entrée grammaticale constitue une barrière cognitive souvent
infranchissable, renforçant l'idée que « ce n'est pas pour eux ».\cite{I-3-1-ref53}



\subsection{Culturalisation vs rigueur linguistique}

Des tentatives de modernisation pédagogique ont eu lieu, cherchant à privilégier l'approche
culturelle, la mythologie, l'étymologie, pour rendre la matière plus accessible et
attractive.\cite{I-3-1-ref50} Cependant, ces tentatives se heurtent souvent à la résistance d'une
partie du corps professoral et des parents d'élèves conservateurs, pour qui la valeur du grec réside
précisément dans sa difficulté. Si le grec devient « facile » ou « ludique », il perd sa fonction de
signal de compétence et de distinction. Le débat pédagogique est donc indissociable du débat social
sur la fonction de l'option.\cite{I-3-1-ref56}



\subsection{Conclusion}

Au terme de cette analyse approfondie, il apparaît que le grec ancien a joué un rôle central,
quoique discret, dans l'adaptation des classes sociales favorisées au « choc » de la réforme Haby.
Loin de disparaître dans le creuset du collège unique, il s'est transformé en une ressource
stratégique rare, une « option-refuge » permettant de contourner la sectorisation (carte scolaire)
et de recréer des filières d'élite au sein des établissements (classes de niveau).

Ce phénomène illustre parfaitement le concept de « démocratisation ségrégative » de Pierre Merle.
L'ouverture quantitative du collège s'est accompagnée d'une fermeture sociale des filières
distinctives. Les dispositifs administratifs comme le motif de dérogation pour « parcours scolaire
particulier » ont institutionnalisé cette faille, permettant aux initiés de transformer un choix
pédagogique en stratégie résidentielle et sociale.

Ainsi, cinquante ans après la réforme Haby, le grec ancien demeure un bastion. Sa survie précaire
dans le système public ne tient pas seulement à l'attachement patrimonial à la culture antique, mais
aussi à son utilité sociologique majeure : garantir à une minorité d'élèves un parcours protégé, une
scolarité entre soi, et la certification d'une excellence scolaire héritée. L'échec du collège
unique à réaliser la mixité sociale se lit, en creux, dans la persistance de ces forteresses
hellénistes.

